\section{Metodologia}

O projeto será desenvolvido em etapas incrementais e iterativas, garantindo o rigor metodológico característico da otimização contínua.

\subsection{Estudo Dirigido e Revisão Bibliográfica}
Inicialmente, será conduzido um estudo aprofundado dos métodos de região de confiança para otimização não-suave \cite{leconte2023complexity}, comportamento algorítmico não monótono e métodos de gradiente proximal indefinido \cite{Leconte_2024} desenvolvidos pelo grupo na Polytechnique Montréal. Paralelamente, haverá um treinamento focado na arquitetura do \texttt{JuliaSmoothOptimizers} para garantir que os desenvolvimentos posteriores estejam alinhados com as boas práticas de engenharia de software da comunidade.

\subsection{Desenvolvimento de Algoritmos Híbridos Ativos}
A metodologia de pesquisa focará na concepção de um modelo algorítmico híbrido. O componente desenvolvido no Brasil (NSPG + buscas combinatórias) atuará nas iterações iniciais do processo de otimização, promovendo esparsidade e identificando rapidamente um conjunto ativo (as variáveis não nulas). 

Adicionalmente, investigaremos o uso de matrizes de segunda ordem estruturadas no subproblema do operador proximal (operadores proximais de métrica variável). Embora abordagens recentes utilizando métricas diagonais indefinidas simples tenham apresentado resultados inferiores em comparação aos métodos espectrais padrão \cite{Leconte_2024}, a adoção de estruturas levemente mais densas, como matrizes tridiagonais, pode proporcionar um ganho significativo de convergência sobre o modelo espectral (múltiplo da identidade). De fato, nossos experimentos preliminares em mínimos quadrados com regularização L0 e estruturas tridiagonais já mostraram resultados promissores, demonstrando que tais subproblemas mantêm a rapidez computacional ao mesmo tempo em que oferecem alta precisão na recuperação do suporte.

Uma vez identificado o conjunto ativo pelos pré-processadores globais, o algoritmo aproveitará a descida de coordenadas e as buscas combinatórias para refinar estritamente a solução sobre as variáveis não nulas. Essa estratégia de \textit{active-set} reduzirá drasticamente a complexidade computacional, combinando a rápida navegação do espaço por operadores proximais estruturados com a convergência local precisa e isolamento de suporte dos refinadores CD.

\subsection{Avaliação Numérica e Benchmarking}
A implementação será feita na linguagem Julia, devido ao seu alto desempenho, que se alinha com o escopo do projeto no Brasil e a infraestrutura JSO no Canadá. Serão conduzidos experimentos extensivos comparando o método híbrido proposto contra algoritmos do estado da arte (como PGCCD e o pacote L0Learn) em instâncias sintéticas de alta dimensão e conjuntos de dados reais de aprendizado estatístico. As métricas de avaliação incluirão tempo de CPU, precisão de suporte, escalabilidade e qualidade da solução final em termos da função objetivo restrita.
